\documentclass[12pt,a4paper]{article}
\usepackage[utf8]{inputenc}
\usepackage[T1]{fontenc}
\usepackage[english]{babel}
\usepackage[margin=2.5cm]{geometry}
\usepackage{amsmath,amssymb,amsthm}
\usepackage{booktabs}
\usepackage{array}
\usepackage{xcolor}
\usepackage{hyperref}
\usepackage{titlesec}
\usepackage{fancyhdr}
\usepackage{enumitem}
\usepackage{listings}
\usepackage{float}

\usepackage{parskip}

\definecolor{primary}{RGB}{44,62,80}
\definecolor{accent}{RGB}{52,152,219}
\definecolor{codebg}{RGB}{245,245,245}
\definecolor{codeframe}{RGB}{200,200,200}

\hypersetup{colorlinks=true,linkcolor=primary,urlcolor=accent,citecolor=accent,
  pdftitle={Monte Carlo Simulation and Statistical Modeling},pdfauthor={Abdellah Kahlaoui}}

\titleformat{\section}{\large\bfseries\color{primary}}{{\thesection}}{1em}{}[\color{accent}\hrule\vspace{2pt}]
\titleformat{\subsection}{\normalsize\bfseries\color{primary}}{\thesubsection}{1em}{}

\lstset{language=Python,basicstyle=\ttfamily\small,backgroundcolor=\color{codebg},
  frame=single,rulecolor=\color{codeframe},breaklines=true,
  keywordstyle=\color{accent}\bfseries,commentstyle=\color{gray}\itshape,
  stringstyle=\color{orange!80!black},numbers=left,numberstyle=\tiny\color{gray},
  numbersep=8pt,showstringspaces=false,captionpos=b}

\pagestyle{fancy}\fancyhf{}
\lhead{\small\color{gray}Monte Carlo Simulation \& Statistical Modeling}
\rhead{\small\color{gray}Abdellah Kahlaoui}
\cfoot{\thepage}
\renewcommand{\headrulewidth}{0.4pt}

\newtheorem{theorem}{Theorem}
\newtheorem{definition}{Definition}
\DeclareMathOperator{\E}{\mathbb{E}}
\DeclareMathOperator{\Var}{Var}

\begin{document}

\begin{titlepage}
  \centering\vspace*{2cm}
  {\color{primary}\rule{\linewidth}{1.5pt}}\\[0.4cm]
  {\LARGE\bfseries\color{primary}Monte Carlo Simulation\\[0.3cm]and Statistical Modeling}\\[0.4cm]
  {\color{primary}\rule{\linewidth}{1.5pt}}\\[1.5cm]
  {\large\itshape Analysis of Complex Probability Distributions,\\
    CLT Convergence Verification,\\and European Option Pricing}\\[2cm]
  {\large\textbf{Abdellah Kahlaoui}}\\[0.3cm]
  {\normalsize Master of Applied Mathematics}\\
  {\normalsize FST Settat, Morocco}\\[0.5cm]
  {\normalsize\href{mailto:kahlaouiabdellah6@gmail.com}{kahlaouiabdellah6@gmail.com}}\\
  {\normalsize\href{https://github.com/AbdellahKah}{github.com/AbdellahKah}}\\[2cm]
  {\normalsize November 2025 -- January 2026}\\[1.5cm]
  \vfill
  {\small\color{gray}Source code: \href{https://github.com/AbdellahKah/monte-carlo-statistical-modeling}{github.com/AbdellahKah/monte-carlo-statistical-modeling}}
\end{titlepage}

\begin{abstract}
\noindent
This report presents a comprehensive Monte Carlo simulation study of four canonical probability
distributions: Normal, Log-Normal, Exponential, and Poisson. We develop a vectorized simulation
engine in Python capable of generating $10^6$ independent paths per distribution in under one
second. Empirical moments are validated against closed-form theoretical values, and goodness-of-fit
is assessed via Kolmogorov-Smirnov and chi-squared tests. We experimentally confirm the convergence
rate predicted by the Central Limit Theorem, measuring log-log slopes between $-0.487$ and $-0.509$
across all distributions. Finally, we apply the Monte Carlo method to the pricing of a European call
option under Black-Scholes assumptions, achieving a price estimate of $10.4284$ against the
analytical value of $10.4506$, a relative error below $0.3\%$ at $n=10^6$ paths.
\end{abstract}

\tableofcontents\newpage

\section{Introduction}

The Monte Carlo method is a broad class of computational algorithms that rely on repeated random
sampling to obtain numerical estimates of quantities that may be difficult or impossible to compute
analytically. First formalised by Ulam and von Neumann in the 1940s, Monte Carlo methods are
now ubiquitous across quantitative finance, statistical physics, and machine learning.

In quantitative finance, Monte Carlo simulation plays a central role in pricing path-dependent
derivatives, computing Value-at-Risk and Expected Shortfall, calibrating stochastic volatility
models, and simulating credit default scenarios.

The fundamental principle is straightforward: if $X$ is a random variable with distribution $F$,
and we wish to estimate $\mu = \E[g(X)]$, we generate $n$ independent samples $X_1,\ldots,X_n\sim F$
and approximate:
\begin{equation}
  \hat{\mu}_n = \frac{1}{n}\sum_{i=1}^{n}g(X_i) \xrightarrow{\text{a.s.}} \mu
  \quad\text{as }n\to\infty \label{eq:mc}
\end{equation}
by the Strong Law of Large Numbers. The accuracy decays as $\sigma/\sqrt{n}$ by the CLT ---
a rate we verify experimentally in Section~\ref{sec:convergence}.

\section{Theoretical Background}

\subsection{Probability Distributions}

\begin{definition}
$X\sim\mathcal{N}(\mu,\sigma^2)$ has PDF:
$f(x)=\frac{1}{\sigma\sqrt{2\pi}}\exp\!\left(-\frac{(x-\mu)^2}{2\sigma^2}\right)$.
Moments: $\E[X]=\mu$, $\Var(X)=\sigma^2$, skewness $=0$.
\end{definition}

\begin{definition}
$X\sim\text{LogNormal}(\mu,\sigma^2)$ if $\log X\sim\mathcal{N}(\mu,\sigma^2)$.
Moments: $\E[X]=e^{\mu+\sigma^2/2}$, $\Var(X)=(e^{\sigma^2}-1)e^{2\mu+\sigma^2}$.
The log-normal is the canonical model for asset prices in the Black-Scholes framework.
\end{definition}

\begin{definition}
$X\sim\text{Exp}(\lambda)$ has PDF $f(x)=\lambda e^{-\lambda x}$, $x\geq 0$.
Moments: $\E[X]=1/\lambda$, skewness $=2$.
Models inter-arrival times and default intensity in credit risk.
\end{definition}

\begin{definition}
$X\sim\text{Poisson}(\lambda)$: $P(X=k)=\lambda^k e^{-\lambda}/k!$.
Moments: $\E[X]=\Var(X)=\lambda$, skewness $=1/\sqrt{\lambda}$.
Models jump counts in Merton-type jump-diffusion processes.
\end{definition}

\subsection{The Central Limit Theorem}

\begin{theorem}[CLT]
Let $X_1,X_2,\ldots$ be i.i.d.\ with mean $\mu$ and finite variance $\sigma^2>0$. Then:
\[
  \sqrt{n}\,\frac{\bar{X}_n-\mu}{\sigma}\xrightarrow{d}\mathcal{N}(0,1)\quad\text{as }n\to\infty
\]
\end{theorem}
A direct consequence is the Monte Carlo standard error:
\begin{equation}
  SE(\hat{\mu}_n)=\frac{\sigma}{\sqrt{n}}=O(n^{-1/2}) \label{eq:se}
\end{equation}
In log-log space: $\log SE = \log\sigma - \frac{1}{2}\log n$, so the slope is exactly $-0.5$.

\subsection{Goodness-of-Fit Testing}

For continuous distributions we use the \textbf{Kolmogorov-Smirnov} test:
$D_n=\sup_x|\hat{F}_n(x)-F_0(x)|$, where $\hat{F}_n$ is the empirical CDF.

For the discrete Poisson distribution we use \textbf{Pearson's chi-squared} test:
$\chi^2=\sum_k(O_k-E_k)^2/E_k$, with bins merged to ensure $E_k\geq 5$.

\section{Implementation}

\subsection{Simulation Engine}

The engine uses \texttt{numpy.random.default\_rng} (PCG64 algorithm), generating all $10^6$
paths in a single vectorized call:

\begin{lstlisting}[caption={Vectorized Normal sampler}]
def simulate_normal(self, mu=0.0, sigma=1.0):
    samples = self.rng.normal(loc=mu, scale=sigma,
                              size=self.n_paths)
    return self._build_result("Normal",
                              {"mu": mu, "sigma": sigma},
                              samples)
\end{lstlisting}

Skewness and kurtosis are computed via vectorized moment operations:
\begin{lstlisting}[caption={Vectorized moment computation}]
centred = samples - mean
skew    = (centred**3).mean() / (std**3 + 1e-12)
kurt    = (centred**4).mean() / (std**4 + 1e-12) - 3.0
\end{lstlisting}

\subsection{Convergence Analyser}

A single batch of $10^6$ samples is drawn once, then sliced at 60 log-spaced checkpoints
from $n=100$ to $n=10^6$. The convergence rate is estimated by OLS regression of
$\log(SE_n)$ on $\log(n)$.

\section{Results --- Distribution Analysis}

\subsection{Empirical vs Theoretical Moments}

\begin{table}[H]
\centering
\caption{Empirical vs theoretical moments --- $n=10^6$ paths}
\renewcommand{\arraystretch}{1.3}
\begin{tabular}{lcccccc}
\toprule
\textbf{Distribution} & \textbf{Emp.\ Mean} & \textbf{Th.\ Mean}
  & \textbf{Emp.\ Std} & \textbf{Th.\ Std}
  & \textbf{Emp.\ Skew} & \textbf{Th.\ Skew}\\
\midrule
Normal ($\mu{=}0,\sigma{=}1$)         & 0.000098 & 0.000000 & 1.000483 & 1.000000 & $-$0.001 & 0.000\\
Log-Normal ($\mu{=}0,\sigma{=}0.5$)   & 1.133568 & 1.133148 & 0.603999 & 0.603901 & 1.755 & 1.750\\
Exponential ($\lambda{=}2$)           & 0.499858 & 0.500000 & 0.498916 & 0.500000 & 1.997 & 2.000\\
Poisson ($\lambda{=}5$)               & 5.000647 & 5.000000 & 2.235591 & 2.236068 & 0.448 & 0.447\\
\bottomrule
\end{tabular}
\end{table}

All empirical moments agree with theoretical values to at least 4 significant figures.

\subsection{Goodness-of-Fit Tests}

\begin{table}[H]
\centering
\caption{Goodness-of-fit results --- subsample of $10^4$ paths}
\renewcommand{\arraystretch}{1.3}
\begin{tabular}{lcccc}
\toprule
\textbf{Distribution} & \textbf{Test} & \textbf{Statistic} & \textbf{p-value} & \textbf{Result}\\
\midrule
Normal      & KS       & 0.010865 & 0.1871 & Pass\\
Log-Normal  & KS       & 0.007083 & 0.6946 & Pass\\
Exponential & KS       & 0.008932 & 0.3999 & Pass\\
Poisson     & $\chi^2$ & 16.421   & 0.2271 & Pass\\
\bottomrule
\end{tabular}
\end{table}

All p-values are well above $\alpha=0.05$, confirming each sampler correctly implements
the target distribution.

\section{Convergence Analysis}
\label{sec:convergence}

\subsection{Experimental Convergence Rates}

\begin{table}[H]
\centering
\caption{Empirical convergence rates vs CLT prediction of $-0.50$}
\renewcommand{\arraystretch}{1.3}
\begin{tabular}{lccc}
\toprule
\textbf{Distribution} & \textbf{True Mean} & \textbf{Empirical Slope} & \textbf{Deviation}\\
\midrule
Normal ($\mu{=}0,\sigma{=}1$)       & 0.0000 & $-$0.5091 & $+$0.0091\\
Log-Normal ($\mu{=}0,\sigma{=}0.5$) & 1.1331 & $-$0.5031 & $+$0.0031\\
Exponential ($\lambda{=}2$)         & 0.5000 & $-$0.4867 & $-$0.0133\\
Poisson ($\lambda{=}5$)             & 5.0000 & $-$0.5066 & $+$0.0066\\
\bottomrule
\end{tabular}
\end{table}

All slopes are within $0.014$ of the theoretical $-0.50$. Small deviations reflect
finite-sample variability in the OLS estimate.

\subsection{Implications for Simulation Design}

From equation~\eqref{eq:se}, halving the standard error requires quadrupling the number of
paths. This motivates variance reduction techniques: antithetic variates (using $Z$ and $-Z$),
control variates, and quasi-Monte Carlo sequences (Sobol, Halton) achieving $O((\log n)^d/n)$.

\section{Application: European Call Option Pricing}

\subsection{Black-Scholes Framework}

Under the risk-neutral measure, the stock price follows:
\[
  S_T = S_0\exp\!\left[\left(r-\tfrac{1}{2}\sigma^2\right)T+\sigma\sqrt{T}\,Z\right],
  \quad Z\sim\mathcal{N}(0,1)
\]
The European call value is:
\[
  C = e^{-rT}\,\E^{\mathbb{Q}}\!\left[\max(S_T-K,\,0)\right]
\]

\subsection{Results}

\begin{table}[H]
\centering
\caption{European call option pricing --- $S_0{=}K{=}100$, $T{=}1$yr, $r{=}5\%$, $\sigma{=}20\%$}
\renewcommand{\arraystretch}{1.3}
\begin{tabular}{lcc}
\toprule
\textbf{Method} & \textbf{Price} & \textbf{Relative Error}\\
\midrule
Black-Scholes analytical  & 10.4506 & ---\\
Monte Carlo ($n=10^6$)    & 10.4284 & $0.21\%$\\
\midrule
\multicolumn{3}{l}{\small 95\% CI: $(10.4035,\;10.4533)$ \quad Std error: $0.01268$}\\
\bottomrule
\end{tabular}
\end{table}

The Black-Scholes analytical price falls within the Monte Carlo 95\% confidence interval,
confirming the statistical validity of the estimate.

\section{Conclusion}

This project has demonstrated the full Monte Carlo pipeline from simulation to financial
application. Key findings:
\begin{enumerate}[noitemsep]
  \item A vectorized Python engine simulates $10^6$ paths per distribution in under one second.
  \item Empirical moments match theory to 4+ significant figures; all goodness-of-fit tests pass.
  \item CLT convergence rates are confirmed experimentally: slopes between $-0.487$ and $-0.509$.
  \item Monte Carlo option pricing achieves $0.21\%$ relative error versus Black-Scholes at $10^6$ paths.
\end{enumerate}
Future work includes variance reduction techniques and extension to path-dependent payoffs.
A separate ongoing project applies neural networks to stochastic volatility model calibration.

\section*{References}
\begin{enumerate}[label={[\arabic*]},noitemsep]
  \item Glasserman, P. (2003). \textit{Monte Carlo Methods in Financial Engineering}. Springer.
  \item Jäckel, P. (2002). \textit{Monte Carlo Methods in Finance}. Wiley.
  \item Black, F. \& Scholes, M. (1973). The Pricing of Options and Corporate Liabilities.
        \textit{Journal of Political Economy}, 81(3), 637--654.
  \item Harris, C.R. et al.\ (2020). Array programming with NumPy. \textit{Nature}, 585, 357--362.
  \item Virtanen, I. et al.\ (2020). SciPy 1.0. \textit{Nature Methods}, 17, 261--272.
\end{enumerate}

\vfill
\begin{center}
  \color{gray}\small\rule{0.4\linewidth}{0.4pt}\\[4pt]
  Abdellah Kahlaoui $\cdot$ FST Settat $\cdot$ 2025--2026\\
  \href{https://github.com/AbdellahKah/monte-carlo-statistical-modeling}{github.com/AbdellahKah/monte-carlo-statistical-modeling}
\end{center}

\end{document}
